\documentclass[carta,final]{lafrox}
\usepackage{simulacion}
\input{datos-proponente}
\usepackage{needspace}% keep headings with their content
\let\SIMoldsection\section
\renewcommand{\section}{\needspace{9\baselineskip}\SIMoldsection}
\begin{document}
\SIMemisor{Tesorería General de la República (emisor simulado)}{Certificado de deuda tributaria}{SIM-TGR-DEUDA-2026-000001}

\noindent Se certifica, para ser presentado en la \LFXlicitacion{} (Art.~39°, numeral~2, de las Bases Administrativas), la situación de deudas tributarias del contribuyente que se individualiza, a la fecha de emisión.

\section*{Contribuyente}
\begin{tablalafrox}{Individualización del contribuyente.}{tab:deu1}{L{5.4cm} Y}
  {\cab{Dato} & \cab{Contenido}}
  \textbf{Rol Único Tributario} & \LFXrut \\
  \textbf{Razón social} & LafroX SpA \\
  \textbf{Domicilio} & \LFXdomicilio \\
  \textbf{Fecha de emisión} & \SIMfecha \\
  \textbf{Fecha de corte de la información} & 22 de noviembre de 2026 \\
\end{tablalafrox}

\section*{Resultado de la consulta}
\noindent A la fecha de corte, el contribuyente \textbf{no registra deudas tributarias morosas}, giros pendientes de pago vencidos ni convenios de pago incumplidos, de aquellos cuya cobranza corresponde a la Tesorería General.

\begin{tablalafrox}{Detalle de formularios y giros revisados.}{tab:deu2}{L{3.2cm} Y L{3.4cm} R{2.6cm}}
  {\cab{Formulario} & \cab{Concepto} & \cab{Períodos revisados} & \cab{Deuda morosa (CLP)}}
  F29 & Declaración mensual y pago simultáneo de IVA y PPM & 11-2021 a 10-2026 & 0 \\
  F50 & Retenciones de impuestos de segunda categoría y adicionales & 11-2021 a 10-2026 & 0 \\
  F22 & Declaración anual de impuesto a la renta & AT 2022 a AT 2026 & 0 \\
  Giros & Giros de impuestos y reajustes pendientes & Todos & 0 \\
  Convenios & Convenios de pago vigentes & Todos & Sin convenios \\
  Contribuciones & Contribuciones de bienes raíces del propietario & Vigentes & Sin deuda \\
  \midrule
  \textbf{Total} & & & \textbf{0} \\
\end{tablalafrox}

\noindent El presente certificado se emite a solicitud del interesado, quien lo presenta para acreditar su situación tributaria ante el CLIENTE. Su vigencia queda limitada a la fecha de emisión.

\vspace{24pt}
\begin{center}
  \rule{6.4cm}{0.5pt}\\[2pt]
  Tesorería General de la República (firma simulada)\\
  {\small\textit{Sin firma electrónica avanzada, sin timbre y sin código de verificación reales}}
\end{center}

\noindent\small Documento recreado con fines académicos. No reemplaza al certificado oficial que emite la Tesorería General de la República.
\end{document}
